\documentclass[12pt]{article}

% Packages
\usepackage[utf8]{inputenc}   % UTF-8 encoding
\usepackage{amsmath, amssymb} % Math symbols
\usepackage{amsthm}           % Theorem environments
\usepackage{geometry}          % Page margins
\usepackage{hyperref}          % For hyperlinks (optional)
\theoremstyle{definition}
\newtheorem{definition}{Definition}[section]

% Page layout
\geometry{a4paper, margin=1in}

\begin{document}

\title{Formal Definition of Hybrid Automaton for Pedestrian Protection}
\author{}
\date{}
\maketitle

A \textbf{hybrid automaton} is formally defined as a tuple
\[
\mathcal{H} = (Q, X, \mathsf{Init}, f, \mathsf{Inv}, E, G, R)
\]
where
\begin{itemize}
    \item $Q$ is a finite set of \textbf{discrete states} (modes or locations).
    \item $X \subseteq \mathbb{R}^n$ is the set of \textbf{continuous variables}.
    \item $\mathsf{Init} \subseteq Q \times X$ is the set of \textbf{initial states}.
    \item $f: Q \times X \to \mathbb{R}^n$ assigns a \textbf{vector field} for continuous evolution in each discrete state: 
    \[
    \dot{x} = f(q, x), \quad q \in Q, \; x \in X
    \]
    \item $\mathsf{Inv}: Q \to 2^X$ assigns an \textbf{invariant} to each discrete state, specifying the set of continuous states where the system can remain in $q$.
    \item $E \subseteq Q \times Q$ is the set of \textbf{edges} representing discrete transitions between states.
    \item $G: E \to 2^X$ assigns a \textbf{guard} to each edge, which is a condition on $x$ that must hold for the transition to occur.
    \item $R: E \times X \to 2^X$ assigns a \textbf{reset map} to each edge, specifying how the continuous state changes during a discrete transition.
\end{itemize}


\section{States}

The automaton has the following discrete states 
\[Q = \{\mathsf{Normal}, \mathsf{SafeWarning}, \mathsf{Throttling}, \mathsf{SoftBraking}, \mathsf{EmergencyBraking} \},\]
each representing a different operational mode of the pedestrian protection system:
\begin{itemize}
    \item Normal: the system operates under normal conditions, driver has full control.
    \item SafeWarning: the system detected a pedestrian potentially crossing (not crossing yet) at a safe-risky TTC. The system issues a warning to the driver.
    \item Throttling: the system detected a pedestrian either crossing at a safe-risky TTC or not crossing yet at risky-critical TTC. The system throttles acceleration.
    \item SoftBraking: the system detected a pedestrian crossing at risky-critical TTC. The system brakes.
    \item EmergencyBraking: the system detected a pedestrian crossing at critical TTC. The system stops the vehicle.
\end{itemize}

\section{Continuous Variables}

The automaton has the following continuous variables:
\[X = \{C_0, \ldots, C_{n-1}, TTC_0, \ldots, TTC_{n-1}, cs_0, \ldots, cs_{n-1}, s_d, s_c, s_u, t\}\]
where:
\begin{itemize}
    \item $C_i \in [0,1]$ is the $i$-th confidence level. For brevity, we denote $C_0, \ldots, C_{n-1}$ as $B_C$;
    \item $TTC_i \in \mathbb{R}_{>0}$ is the $i$-th time-to-collision. For brevity, we denote $TTC_0, \ldots, TTC_{n-1}$ as $B_{TTC}$;
    \item $cs_i \in \{0,1\}$ is the $i$-th crossing state. For brevity, we denote $cs_0, \ldots, cs_{n-1}$ as $B_{cs}$;
    \item $s_d$ is the staleness of the detection data;
    \item $s_c$ is the time since last pedestrian crossing detection;
    \item $s_u$ is the time since uncertain data.
    \item $t$ is the time since last sensor value.
\end{itemize}

\section{Initial States}
The initial state is 
\(\mathsf{Init} = \{(\mathsf{Normal}, X_0)\}\)
where 
\[
\begin{aligned}
    X_0 = &\{C_0 = 0, \ldots, C_{n-1} = 0, TTC_0 = \mathsf{NO\_TTC}, \ldots, TTC_{n-1} = \mathsf{NO\_TTC}, \\
            &cs_0 = 0, \ldots, cs_{n-1} = 0, s_d = 0, s_c = 0, s_u = 0, t = 0\},
\end{aligned}
\]
where $\mathsf{NO\_TTC}$ is a large sentinel value indicating no pedestrian detected.


\section{Thresholds and Helper Functions}

\paragraph{Thresholds.} The automaton uses the following thresholds:
\begin{itemize}
    \item $n$: number of frames;
    \item $\mathsf{TH\_C}$: threshold for detection confidence;
    \item $\mathsf{TH\_D\_stale}$: threshold for detection stale data;
    \item $\mathsf{TH\_C\_stale}$: threshold for crossing stale data;
    \item $\mathsf{MAX\_UNC}$: time upper bound for uncertain camera detection (ms);
    \item $\mathsf{TH\_TTC\_s}$: threshold for safe time-to-collision;
    \item $\mathsf{TH\_TTC\_r}$: threshold for risky time-to-collision;
    \item $\mathsf{TH\_TTC\_c}$: threshold for critical time-to-collision;
    \item $\mathsf{S\_CONS}$: percentage of agreeing frames for safe distance classification; 
    \item $\mathsf{SR\_CONS}$: percentage of agreeing frames for safe-risky distance classification;
    \item $\mathsf{RC\_CONS}$: percentage of agreeing frames for risky-critical distance classification;
    \item $\mathsf{C\_CONS}$: percentage of agreeing frames for critical distance classification;
    \item $\mathsf{CONS}: $ percentage of agreeing frames for classification.
\end{itemize}

\paragraph{Assumptions.} The automaton applies the following assumptions:
\begin{itemize}
    \item $\mathsf{CAM\_FREQ} = 100ms$;
    \item $\mathsf{RT\_H} = 700ms$.
\end{itemize}

\paragraph{Helper functions.} We now define the following helper functions.
\begin{definition}[P helper function]
    The helper function $P$ is defined as follows.
\[
    P(\mathsf{cond}) =
    \begin{cases}
        1 & \text{if } \mathsf{cond} \text{ is true} \\
        0 & \text{otherwise}
    \end{cases}
\]
\end{definition}


\begin{definition}[Sense function]
The sense function $\mathsf{sense}: X \to 2^X$ is defined as follows.
    \[
    \begin{aligned}
        &\mathsf{sense}(\{C_0, \ldots, C_{n-1}, TTC_0, \ldots, TTC_{n-1}, cs_0, \ldots, cs_{n-1}, s_d, s_c, s_u, t\}) = \\
        &\{C, C_1', \ldots, C_{n-1}', TTC, TTC_1', \ldots, TTC_{n-1}', cs, cs_1', \ldots, cs_{n-1}', s_d, s_c, s_u, t'\}
    \end{aligned}
\]
where
\[
    \begin{aligned}
    t' &= 0, \, C_i' = C_{i-1}, \, TTC_i' = TTC_{i-1}, \, cs_i' = cs_{i-1} \text{ with } i = 1, \dots, n-1, \\ 
    C &\in [0,1], \;
    TTC =     
    \begin{cases}
        \mathsf{NO\_TTC} & \text{if } C < \mathsf{TH\_C} \\
        v \in \mathbb{R}_{>0} & \text{otherwise}
    \end{cases} \;
    cs =     
    \begin{cases}
        0 & \text{if } C < \mathsf{TH\_C} \\
        v \in \{0,1\} & \text{otherwise}
    \end{cases} 
    \end{aligned}
\]
\end{definition}



\begin{definition}[Reset timers function]
The reset timer function $\mathsf{reset}: X \to X$ is defined as follows.
    \[
    \begin{aligned}
        &\mathsf{reset}(\{B_C, B_{TTC}, B_{cs},  s_d, s_c, s_u, t\}) = \{B_C, B_{TTC}, B_{cs},  s_d', s_c', s_u',t\}
    \end{aligned}
\]
where 
\[
    \begin{aligned}
    s_d' &=     
    \begin{cases}
        0 & \text{if } \mathsf{detected}(\{B_C, B_{TTC}, B_{cs}, s_d, s_c, s_u, t\}) \\
        s_d & \text{otherwise}
    \end{cases}\\
    s_c' &=     
    \begin{cases}
        0 & \text{if } \mathsf{crossing}(\{B_C, B_{TTC}, B_{cs}, s_d, s_c, s_u, t\}) \\
        s_c & \text{otherwise}
    \end{cases}\\
    s_u' &=     
    \begin{cases}
        0 & \text{if } \mathsf{uncertain}(\{B_C, B_{TTC}, B_{cs}, s_d, s_c, s_u, t\}) \\
        s_u & \text{otherwise}
    \end{cases}\\
    \end{aligned}
\]
\end{definition}

\begin{definition}[Detection]
    The predicate $\mathsf{detected}$ as defined as follows. Recall, $B_C = C_0 \ldots C_{n-1}$.
    \[
        \begin{aligned}
                \mathsf{detected}(\{B_C, B_{TTC}, B_{cs}, s_d, s_c, s_u, t\}) = \sum_{i=0}^{n-1} P(C_i \geq \mathsf{TH\_C}) \geq \lceil \frac{\mathsf{RT\_H}}{ 2 \cdot \mathsf{CAM\_FREQ}} \rceil
        \end{aligned}
    \]
\end{definition}

\begin{definition}[Valid detection]
    The predicate $\mathsf{valid\_d}$ as defined as follows.
    \[
        \begin{aligned}
                \mathsf{valid\_d}(X) = \mathsf{detected}(X) \lor s_d < \mathsf{TH\_D\_stale}
        \end{aligned}
    \]
\end{definition}

\begin{definition}[Safe distance]
The predicate $\mathsf{s\_dist}$ is defined as follows. Recall, $B_{TTC} = TTC_0 \ldots TTC_{n-1}$.
\[
    \begin{aligned}
            \mathsf{s\_dist}&(\{B_C, B_{TTC}, B_{cs}, s_d, s_c, s_u, t\}) = \sum_{i=0}^{k} P(TTC_i > \mathsf{TH\_TTC\_s}) \geq \lceil k \cdot \mathsf{ CONS}\rceil
    \end{aligned}
\]
where $k = \lceil n \cdot \mathsf{S\_CONS} \rceil$.
    
\end{definition}

\begin{definition}[Safe-Risky distance]
The predicate $\mathsf{s\_r\_dist}$ is defined as follows. Recall, $B_{TTC} = TTC_0 \ldots TTC_{n-1}$.
\[
    \begin{aligned}
        \mathsf{s\_r\_dist}&(\{B_C, B_{TTC}, B_{cs}, s_d, s_c, s_u, t\}) = \\
        &\sum_{i=0}^{k} P(\mathsf{TH\_TTC\_r} < TTC_i <= \mathsf{TH\_TTC\_s}) \geq \lceil k \cdot \mathsf{CONS}\rceil
    \end{aligned}
\]
where $k = \lceil n \cdot \mathsf{SR\_CONS} \rceil$.
    
\end{definition}

\begin{definition}[Risky-Critical distance]
The predicate $\mathsf{r\_c\_dist}$ is defined as follows. Recall, $B_{TTC} = TTC_0 \ldots TTC_{n-1}$.
\[
    \begin{aligned}
        \mathsf{r\_c\_dist}&(\{B_C, B_{TTC}, B_{cs}, s_d, s_c, s_u, t\}) = \\ 
        &\sum_{i=0}^{k} P(\mathsf{TH\_TTC\_c} < TTC_i <= \mathsf{TH\_TTC\_r}) \geq \lceil k \cdot \mathsf{CONS}\rceil
    \end{aligned}
\]
where $k = \lceil n \cdot \mathsf{RC\_CONS} \rceil$.    
\end{definition}

\begin{definition}[Critical distance]
The predicate $\mathsf{c\_dist}$ is defined as follows. Recall, $B_{TTC} = TTC_0 \ldots TTC_{n-1}$.
\[
    \begin{aligned}
        \mathsf{c\_dist}&(\{B_C, B_{TTC}, B_{cs}, s_d, s_c, s_u, t\}) = \sum_{i=0}^{k} P(TTC_i \leq \mathsf{TH\_TTC\_c}) \geq \lceil k \cdot \mathsf{CONS}\rceil
    \end{aligned}
\]
where $k = \lceil n \cdot \mathsf{C\_CONS} \rceil$.    
\end{definition}

\begin{definition}[Crossing]
    The predicate $\mathsf{crossing}$ as defined as follows. Recall, $B_{cs} = cs_0 \ldots cs_{n-1}$.
    \[
        \begin{aligned}
            \mathsf{crossing}(\{B_C, B_{TTC}, B_{cs}, s_d, s_c, s_u, t\}) = \sum_{i=0}^{n-1} P(cs_i = 1) \geq \lceil \frac{\mathsf{RT\_H}}{ 2 \cdot \mathsf{CAM\_FREQ}} \rceil
        \end{aligned}
    \]
\end{definition}

\begin{definition}[Valid crossing]
    The predicate $\mathsf{valid\_c}$ as defined as follows.
    \[
        \mathsf{valid\_c}(X) = \mathsf{crossing}(X) \lor s_c < \mathsf{TH\_C\_stale}
    \]
\end{definition}

\begin{definition}[Uncertain data]
    The predicate $\mathsf{uncertain}$ as defined as follows.
    \[
        \mathsf{uncertain}(X) = \neg \mathsf{s\_dist}(X) \land \neg \mathsf{s\_r\_dist}(X) \land \neg \mathsf{r\_c\_dist}(X) \land \neg \mathsf{c\_dist}(X)
    \]
\end{definition}

\begin{definition}[Valid uncertain]
    The predicate $\mathsf{uncertain}$ as defined as follows.
    \[
        \mathsf{valid\_u}(X) = \mathsf{uncertain}(X) \land s_u < \mathsf{MAX\_UNC}
    \]
\end{definition}


\section{Vector Field}
The vector field $f$ defines the continuous evolution of the variables in each discrete state: 
\[
\dot{\mathsf{sns}} = 0, \quad \dot{C}_i = 0, \quad \dot{TTC}_i = 0, \quad \dot{cs}_i = 0, \quad \forall i \in [0,n-1], 
\]

\[
\dot{s}_d = 1000, \quad \dot{s}_c = 1000, \quad \dot{s}_u = 1000, \quad \dot{t} = 1000
\]

Equivalently, the vector field can be written as:
\[
X = 
\begin{bmatrix}
    C_0 \\ \vdots \\ C_{n-1} \\ 
    TTC_0 \\ \vdots \\ TTC_{n-1} \\ 
    cs_0 \\ \vdots \\ cs_{n-1} \\ 
    s_d \\ s_c \\ s_u \\ t
\end{bmatrix}, \quad
f(q, X) =
\begin{bmatrix}
    0 \\ \vdots \\ 0 \\ 
    0 \\ \vdots \\ 0 \\ 
    0 \\ \vdots \\ 0 \\ 0 \\
    1000 \\ 1000 \\ 1000 \\ 1000
\end{bmatrix}, \quad \forall q \in Q
\]

\section{Invariants}

The invariants $\mathsf{Inv}: Q \to 2^X$ specify conditions under which the automaton can stay in each discrete state.
Notice that $X$ is a shorthand for 
\[\{C_0, \ldots, C_{n-1}, TTC_0, \ldots, TTC_{n-1}, cs_0, \ldots, cs_{n-1}, s_d, s_c, s_u, t\}\]

\[
\begin{aligned}
&\mathsf{Inv}(\mathsf{Normal}) = \{ X \mid (\neg \mathsf{valid\_d}(X) \lor \mathsf{s\_dist}(X) \lor \mathsf{uncertain}(X)) \land t < \mathsf{CAM\_FREQ} \} \\
&\mathsf{Inv}(\mathsf{SafeWarning}) = \{ X \mid \big((\mathsf{valid\_d}(X) \land \neg \mathsf{valid\_c}(X)) \lor \mathsf{valid\_u}(X)\big) \land t < \mathsf{CAM\_FREQ}\} \\
&\mathsf{Inv}(\mathsf{Throttling}) = \{ X \mid \mathsf{valid\_d}(X) \land \mathsf{valid\_c}(X) \land (\mathsf{s\_r\_dist}(X) \lor  \mathsf{valid\_u}(X) ) \land t < \mathsf{CAM\_FREQ} \} \\
&\mathsf{Inv}(\mathsf{SoftBraking}) = \{ X \mid \mathsf{valid\_d}(X) \land \mathsf{valid\_c}(X) \land (\mathsf{r\_c\_dist}(X) \lor  \mathsf{valid\_u}(X) ) \land t < \mathsf{CAM\_FREQ} \} \\
&\mathsf{Inv}(\mathsf{EmergencyBraking}) = \{ X \mid \mathsf{valid\_c}(X) \land t < \mathsf{CAM\_FREQ}\}
\end{aligned}
\]

\section{Edges, Guards, and Resets}

\subsection{Edges}

The set of edges $E \subseteq Q \times Q$ is defined as $E = \{ e_1, e_2, \dots, e_{22} \}$, where
\[
\begin{aligned}
e_1 &= (\mathsf{Normal}, \mathsf{Normal}) \\
e_2 &= (\mathsf{Normal}, \mathsf{SafeWarning}) \\
e_3 &= (\mathsf{Normal}, \mathsf{Throttling}) \\
e_4 &= (\mathsf{Normal}, \mathsf{SoftBraking}) \\
e_5 &= (\mathsf{Normal}, \mathsf{EmergencyBraking}) \\
e_6 &= (\mathsf{SafeWarning}, \mathsf{Normal}) \\
e_7 &= (\mathsf{SafeWarning}, \mathsf{SafeWarning}) \\
e_8 &= (\mathsf{SafeWarning}, \mathsf{Throttling}) \\
e_9 &= (\mathsf{SafeWarning}, \mathsf{SoftBraking}) \\
e_{10} &= (\mathsf{SafeWarning}, \mathsf{EmergencyBraking}) \\
e_{11} &= (\mathsf{Throttling}, \mathsf{Normal}) \\
e_{12} &= (\mathsf{Throttling}, \mathsf{SafeWarning}) \\
e_{13} &= (\mathsf{Throttling}, \mathsf{Throttling}) \\
e_{14} &= (\mathsf{Throttling}, \mathsf{SoftBraking}) \\
e_{15} &= (\mathsf{Throttling}, \mathsf{EmergencyBraking}) \\
e_{16} &= (\mathsf{SoftBraking}, \mathsf{Normal}) \\
e_{17} &= (\mathsf{SoftBraking}, \mathsf{SafeWarning}) \\
e_{18} &= (\mathsf{SoftBraking}, \mathsf{Throttling}) \\
e_{19} &= (\mathsf{SoftBraking}, \mathsf{SoftBraking}) \\
e_{20} &= (\mathsf{SoftBraking}, \mathsf{EmergencyBraking}) \\
e_{21} &= (\mathsf{EmergencyBraking}, \mathsf{Normal}) \\
e_{22} &= (\mathsf{EmergencyBraking}, \mathsf{EmergencyBraking})
\end{aligned}
\]

\subsection{Guards}

The guard conditions $G: E \to 2^{X}$ for each edge are defined as follows. Notice that $X$ is a shorthand for 
\[\{C_0, \ldots, C_{n-1}, TTC_0, \ldots, TTC_{n-1}, cs_0, \ldots, cs_{n-1}, s_d, s_c, s_u, t\}\]

\[
\begin{aligned}
G(e_1)  &= \{X \mid t \ge \mathsf{CAM\_FREQ}\} \\
G(e_2)  &= \{X \mid \mathsf{valid\_d}(X) \land \neg \mathsf{valid\_c}(X) \land  t < \mathsf{CAM\_FREQ} \} \\
G(e_3)  &= \{X \mid \mathsf{valid\_d}(X) \land \mathsf{valid\_c}(X) \land \mathsf{s\_r\_dist}(X) \land t < \mathsf{CAM\_FREQ} \} \\
G(e_4)  &= \{X \mid \mathsf{valid\_d}(X) \land \mathsf{valid\_c}(X) \land \mathsf{r\_c\_dist}(X) \land t < \mathsf{CAM\_FREQ}\} \\
G(e_5)  &= \{X \mid \mathsf{valid\_d}(X) \land \mathsf{valid\_c}(X) \land \mathsf{c\_dist}(X) \land t < \mathsf{CAM\_FREQ}\} \\
G(e_6)  &= \{X \mid (\neg \mathsf{valid\_d}(X) \lor \mathsf{s\_dist}(X) \lor \neg \mathsf{valid\_u}(X)) \land  t < \mathsf{CAM\_FREQ} \} \\
G(e_7)  &= G(e_1) \\
G(e_8) &= G(e_3) \\ 
\end{aligned}
\]
\[
\begin{aligned}
G(e_9) &= G(e_4) \\
G(e_{10}) &= G(e_5) \\ 
G(e_{11}) &= G(e_6) \\
G(e_{12}) &= G(e_2) \\ 
G(e_{13}) &= G(e_1) \\
G(e_{14}) &= G(e_4) \\
G(e_{15}) &= G(e_5) \\
G(e_{16}) &= G(e_6) \\
G(e_{17}) &= G(e_2) \\
G(e_{18}) &= G(e_3) \\
G(e_{19}) &= G(e_1) \\
G(e_{20}) &= G(e_5)\\
G(e_{21}) &= \{ X \mid \neg \mathsf{valid\_c}(X)\} \\
G(e_{22}) &= G(e_1) 
\end{aligned}
\]

\subsection{Resets}

Resets $ R(e,X) = h(X), \forall e \in E$, where 

\[
    h(X) = 
    \begin{cases}
        \mathsf{sense}(X) & \text{if } t \ge \mathsf{CAM\_FREQ} \\
        \{\mathsf{reset}(X)\} & \text{otherwise}
    \end{cases}
\]
Notice that $X$ is a shorthand for 
\[\{C_0, \ldots, C_{n-1}, TTC_0, \ldots, TTC_{n-1}, cs_0, \ldots, cs_{n-1}, s_d, s_c, s_u, t\}\]





\end{document}
